\section{Proof-Class Barrier Taxonomy}

This section records rigorous results showing that several natural proof
programs cannot close the sharp-rate gap.  Unlike the structural negatives of
\Cref{sec:structural-negatives}, the emphasis here is methodological: each
subsection isolates a named class of arguments and proves that the class is
insufficient.  We also keep the strategy-dependence labels explicit.  In
particular, the headline upper bound of \Cref{thm:main-upper-bound} is proved
for the specific Shortener strategy $\sigma_{15}$ and does not use the
max-unresolved-harmonic-degree strategy $\sigma^*$.

\subsection{Transversal-Integrality and Sherali--Adams}

The first barrier is purely static.  It shows that low-level fractional or
Sherali--Adams relaxations cannot certify the large transversals needed for the
late separator program.

% source: researcher-58-pro-R52-integrality-barrier-audit-confirmed-sharpened.md
\begin{theorem}[Transversal-integrality barrier]\label{thm:sa-barrier}
Fix $0<\alpha<1$.  Let $P$ be an $N$-point set, put
\[
\ell=\lfloor \log N\rfloor,
\qquad
q=\lfloor \alpha N\rfloor,
\]
and write $\tau_{\mathbb Z}(\mathcal C)$ and $\tau_f(\mathcal C)$ for the
integral and fractional transversal numbers of an $\ell$-uniform family
$\mathcal C\subseteq \binom{P}{\ell}$.  Then for all sufficiently large $N$
there exists such a family $\mathcal C$ with
\[
\tau_{\mathbb Z}(\mathcal C)> \alpha N,
\qquad
\tau_f(\mathcal C)\le \frac{N}{\ell}.
\]
Moreover, for every $\ell$-uniform family $\mathcal H\subseteq \binom{P}{\ell}$
and every $0\le r<\ell$,
\[
\mathrm{SA}_r(\mathcal H)\le \frac{N}{\ell-r}.
\]
In particular, if $r\le \ell/2$, then
\[
\mathrm{SA}_r(\mathcal C)\le \frac{2N}{\ell}=(2+o(1))\frac{N}{\log N},
\]
while $\tau_{\mathbb Z}(\mathcal C)>\alpha N$, so the level-$r$ integrality
gap is at least $(\alpha/2+o(1))\log N$.
\end{theorem}

\begin{proof}[Proof sketch]
The family $\mathcal C$ is a random $q$-avoiding covering design.  One chooses
\[
s\asymp \delta^{-1}\log \binom{N}{q},
\qquad
\delta=\frac{\binom{N-q}{\ell}}{\binom{N}{\ell}},
\]
independent $\ell$-subsets of $P$.  The union bound shows that, with positive
probability, every $q$-subset of $P$ misses at least one chosen block; this is
equivalent to $\tau_{\mathbb Z}(\mathcal C)>q\ge \alpha N+O(1)$.  On the other
hand, the uniform fractional weight $\ell/N$ on each vertex gives
$\tau_f(\mathcal C)\le N/\ell$.

The Sherali--Adams estimate is cleaner and fully deterministic.  For an
$\ell$-uniform family $\mathcal H$, the standard symmetric level-$r$ assignment
with marginal $(\ell-r)/N$ is feasible, yielding
$\mathrm{SA}_r(\mathcal H)\le N/(\ell-r)$.  The resulting gap is therefore
already logarithmic at all levels $r\le \ell/2$.
\end{proof}

\begin{remark}
This theorem is strategy-independent: it is a static obstruction to any proof
that hopes to derive a strong game upper bound from low-level separator LPs or
their Sherali--Adams tightenings.
\end{remark}

\subsection{Fixed-Rank $q$-Shadow / Covering Dichotomy}

The second barrier isolates the precise point where the spectral packet program
can fail.  Either there is an honest legal separator of the expected size, or
Shortener's own previously created $q$-shadow already occupies almost all of
the separator layer.

% source: researcher-59-pro-R53-q-shadow-dichotomy-audit-confirmed-sharpened.md
\begin{theorem}[Fixed-rank $q$-shadow dichotomy]\label{thm:q-shadow-dichotomy}
Let $P$ be a set of size $K=h+L$ with $h,L\ge 1$, and fix $1\le q\le h$.  Set
\[
Y=\binom{P}{L},
\qquad
Q_q=\binom{P}{q}.
\]
For a blocker family $\mathcal D\subseteq 2^P$ and a played complement family
$\mathcal C\subseteq Y$, let $F_q(\mathcal D)\subseteq Q_q$ be the
comparability shadow of $\mathcal D$ on the $q$-layer, let
$B_q(C)=\{Q\in Q_q:Q\cap C=\varnothing\}$, and define the legal separator set
\[
A_q(\mathcal D,\mathcal C)
:=
Q_q\setminus \Bigl(F_q(\mathcal D)\cup \bigcup_{C\in\mathcal C} B_q(C)\Bigr).
\]
If $R\subseteq Y$ has density
\[
r=\frac{|R|}{|Y|}>0
\]
and no legal $q$-separator in $A_q(\mathcal D,\mathcal C)$ captures at least
$\frac12\delta_q |R|$ live complements, where
\[
\delta_q=\frac{\binom{K-q}{L}}{\binom{K}{L}},
\qquad
\lambda_q^2=\frac{h(K-q)}{qL},
\]
then
\[
\sigma_q(\mathcal D)+|\mathcal C|\delta_q
>
1-\frac{4\lambda_q^2}{r},
\qquad
\sigma_q(\mathcal D):=\frac{|F_q(\mathcal D)|}{|Q_q|}.
\]
At the central scale $q\sim 2(\log h)^2$, $L\sim h/\log h$, one has
\[
\delta_q=(e+o(1))h^{-2},
\qquad
\lambda_q^2=(2+o(1))\frac{\log h}{h},
\]
so if $\sigma_q(\mathcal D)=o(1)$, then necessarily
\[
|\mathcal C|>\left(\frac{1}{e}-o(1)\right)h^2.
\]
\end{theorem}

\begin{proof}[Proof sketch]
The proof packages the Johnson-scheme separator theorem in the exact form
needed for the packet game.  If no legal $q$-separator captures half of the
expected mass, then the live set $R$ has negligible projection onto the legal
separator subspace; quantitatively this forces
$|A_q(\mathcal D,\mathcal C)|/|Q_q|<4\lambda_q^2/r$.  Rewriting the complement
of $A_q$ gives the displayed bound on
$\sigma_q(\mathcal D)+|\mathcal C|\delta_q$.

The central-scale constants are then obtained by direct asymptotics for
$\delta_q$ and $\lambda_q^2$.  The upshot is conceptual: spectral starvation is
possible only after Shortener has already accumulated a macroscopic amount of
her own $q$-shadow or has played on the order of $h^2$ complements.
\end{proof}

\subsection{The Separator-Only Barrier}

We next record the cleanest theorem to emerge from the late
separator-first program: any proof that uses only finitely many odd
lower-half separators, with no genuine composite fallback, cannot yield even an
$o(n)$ upper bound.

% source: phase4/r56_separator_only_barrier_note.md
\begin{theorem}[Finite odd-carrier separator-only barrier]\label{thm:separator-only}
Let $r_1(n)=n(\log\log n)^2/\log n$.  No finite odd-carrier
separator-only closure can prove $L(n)=O(r_1(n))$.  More precisely, if
$S_n\subseteq [2,n/2]$ is an odd separator family with
\[
|S_n|=o(r_1(n)),
\qquad
\sum_{p\mid s\text{ for some }s\in S_n}\frac1p=o(1),
\]
and if the only lower-half composite moves used by the method are elements of
$S_n$, then Prolonger has a legal prefix of length $o(r_1(n))$ after which all
separators in $S_n$ are dead and at least
\[
\frac n2-o(n)
\]
elements of $\U$ remain legal.
\end{theorem}

\begin{proof}[Proof sketch]
If $s\in S_n$ is still legal, Prolonger plays the dyadic multiple
$x=2^a s\in (n/2,n]$, choosing $a$ so that $n/2<x\le n$.  Since all elements of
$S_n$ are odd, any previous lower-half separator comparable to $x$ would have
been comparable to $s$ already, contradicting the legality of $s$.  Distinct
upper-half moves are automatically incomparable.  Thus each legal separator can
be preempted at unit cost.

After at most $2|S_n|=o(r_1(n))$ turns, every separator is dead.  The remaining
illegal elements of $\U$ are those already played in the upper half, together
with the multiples of the prime carrier set of $S_n$; the reciprocal carrier
mass assumption bounds the latter by $o(n)$ via the union bound.  Hence
$|\U_{\mathrm{legal}}|=n/2-o(n)$ at the end of the prefix.
\end{proof}

\begin{remark}
The theorem is a class-level statement.  It does not say that every separator
idea fails; it says that a proof using only finitely many odd separators and no
composite fallback cannot close the game.  This is the correct reading of the
failure of the separator-first route and of the specific strategy
$\tau_{\mathrm{SF}}$.
\end{remark}

\subsection{Strategy-Dependence Classification}

The previous subsections are all strategy-independent.  The next result is a
meta-theorem describing which later ingredients survive strategy replacement
and which do not.

% source: researcher-54-pro-sigma-star-audit-class-C-named-shortener-class.md
\begin{proposition}[Strategy-dependence classification]\label{prop:strategy-dependence}
The rigorous ingredients of the program partition into three classes.
\begin{enumerate}
  \item \emph{Strategy-independent} results: Shield Reduction, Theorem~A, T2,
  the Round~15 $\sigma_{15}$ upper bound, the separator theorems of
  \Cref{thm:sa-barrier,thm:q-shadow-dichotomy,thm:separator-only}, and the
  structural negatives of \Cref{sec:structural-negatives}.
  \item \emph{$\sigma^*$-dependent} results: the R35 state inequality, online
  harmonic domination, and the smallest-legal-prime rule in its unrestricted
  form.
  \item \emph{$\sigma^*$-specific} results: the cooperative-embedding
  construction defeating $\sigma^*$ at cost $n^{o(1)}$, the ST-capture
  refutation in its forced-leaf formulation, and the residual-floor diagnosis
  attached to the $\sigma^*$ bookkeeping framework.
\end{enumerate}
In particular, the load-bearing ingredients of the older packet closure program
do not survive strategy replacement, whereas the headline upper bound of
\Cref{thm:main-upper-bound} is independent of $\sigma^*$.
\end{proposition}

\begin{proof}[Discussion]
This is a classification theorem rather than a new counting argument.  The R54
audit tracks, item by item, which proofs use only legality and combinatorics,
which rely on the max-unresolved-harmonic-degree rule $\sigma^*$, and which are
refutations aimed specifically at the states generated by $\sigma^*$.  The
cooperative-embedding theorem against $\sigma^*$ is decisive here: it shows
that vulnerability of $\sigma^*$ cannot by itself be promoted to a minimax
barrier against all Shortener strategies.
\end{proof}

\subsection{Clearly Labeled Conditional $\sigma^*$ Results}

We record the most useful $\sigma^*$-dependent statements for completeness, but
we emphasize that none of them is used in \Cref{sec:main-upper-bound}.  Their
role is diagnostic, not foundational.

% source: researcher-54-pro-sigma-star-audit-class-C-named-shortener-class.md
% source: phase4/r44_residual_floor_note.md
% source: phase4/r46_st_capture_refutation_note.md
\begin{proposition}[Conditional $\sigma^*$ tools]\label{prop:sigma-star-conditional}
Within the $\sigma^*$ framework one has the following rigorously derived, but
conditional, statements.
\begin{enumerate}
  \item The R35 state inequality and the online harmonic-domination estimate
  control the cumulative useful mass only under the $\sigma^*$ move rule.
  \item The smallest-legal-prime principle is valid only under the additional
  hypothesis that every legal composite has a legal prime factor.
  \item The ST-capture inequality is false in general, even along
  $\sigma^*$-reachable forcing traces.
  \item Even the idealized residual-floor regime with $\mu(F_{\mathrm{useful}})
  =0$ does not by itself reach the conjectural $r_1(n)$ scale.
\end{enumerate}
\end{proposition}

\begin{proof}[References]
Items (1) and the dependence on $\sigma^*$ are catalogued in the R54 audit.
Item (2) is the safe positive lemma isolated in the Round~46 note.  Item (3) is
the ST-capture refutation of Round~46, and item (4) is the residual-floor
diagnosis of Round~44.  Together they explain why the older state-inequality
program is valuable as a diagnostic framework but cannot presently support an
unconditional sharp upper bound.
\end{proof}

\begin{remark}
For reviewers, the key point is simple: the headline theorem of
\Cref{sec:main-upper-bound} is proved for $\sigma_{15}$ and does not use any of
the conditional material in this subsection.
\end{remark}
